\section{Related Work}
\label{section:related-work}
\paragraph{Autoregressive Molecule Generation.} Autoregressive models provide control over the generative process by enabling direct conditioning on prior information, allowing for a more precise and targeted generation of output.
Autoregressive generation has shown success in 2D molecule tasks using SMILE-based methods, as seen in MolMIM~\citep{reidenbach2023improving}, as well as graph-based atom-wise and subgraph-level techniques, as shown in GraphAF \citep{graphaf} and HierVAE \citep{HierVAE}.
Similarly, 3DLinker \citep{3dlinker} and SQUID \citep{squid} showcase the usefulness of 3D autoregressive molecule generation and their ability to leverage conditional information in both atom-wise and subgraph-level settings for 3D linkage and shape-conditioned generative tasks respectively. 
 We note that, unlike prior methods~\citep{squid}, \model does not require a predefined fragment vocabulary.
HERN~\citep{jin2022antibody} further demonstrates the power of hierarchical equivariant autoregressive methods in the task of computational 3D antibody design. Similarly, Pocket2Mol~\citep{pocket2mol} uses autoregressive sampling for structure-based drug design.

\vspace{-1ex}
\paragraph{Protein Docking and Structure-based Drug Design.} Protein docking is a key downstream use case for generating optimal 3D molecule structures. 
Recent research has prominently explored two distinct directions within this field.
The first is blind docking, where the goal is to locate the pocket and generate the optimal ligand to bind~\citep{diffdock}. The second is structure-based drug design (SBDD), where optimal 3D ligands are generated by conditioning on a specific protein pocket. Specifically, the SBDD task focuses on the ability to generate ligands that achieve a low AutoDock Vina score for the CrossDocked2020~\citep{CrossDocked} dataset. AutoDock Vina~\citep{vina} is a widely used molecular docking software that predicts the binding affinity of ligands (drug-like molecules) to target proteins. Autodock Vina takes in the 3D structures of the ligand, target protein, and binding pocket and considers various factors such as van der Waals interactions, electrostatic interactions, and hydrogen bonding between the ligand and target protein to predict the binding affinity. 
We demonstrate how SBDD can be adapted to construct comprehensive MCG benchmarks. In this framework, we evaluate the generative abilities of MCG models by measuring the binding affinities of generated comforters and comparing them to the provided ground truth ligand conformers for a wide array of protein-ligand complexes.

\vspace{-1ex}
\paragraph{SE(3)-Equivariance.} Let $\gX$ and $\gY$ be the input and output vector spaces, respectively, which possess a set of transformations $G$: $G\times\gX\rightarrow\gX$ and $G\times\gY\rightarrow\gY$. The function $\phi: \gX\rightarrow\gY$ is called equivariant with respect to $G$ if, when we apply any transformation to the input, the output also changes via the same transformation or under a certain predictable behavior, \ie,
\begin{definition}
The function $\phi: \gX\mapsto\gY$ is $G$-equivariant if it commutes with any transformation in $G$,
\begin{eqnarray}
\label{eq:equ}
\phi (\rho_{\gX}(g)x) = \rho_{\gY}(g)\phi(x), \forall g\in G,
\end{eqnarray}
where $\rho_{\gX}$ and $\rho_{\gY}$ are the group representations in the input and output space, respectively. Specifically, $\phi$ is called invariant if $\rho_\gY$ is the identity.
\end{definition}
By enforcing SE(3)-equivariance in our probabilistic model,  $p(\mX | \gR)$ remains unchanged for any rototranslation of the approximate conformer $\gR$.
\modelns's architecture is inspired by recent equivariant graph neural network architectures, such as EGNN \citep{EGNN} and PaiNN \citep{painn}, as well as Vector Neuron multi-layer perceptron (VN-MLP) \citep{vn-mlp}.

\section{Loss function}
\label{section:appendix-loss-func}
\vspace{-1ex}
As described in~\sref{sec:methods-algorithmic-details}, \model optimizes the following loss function:
% \vspace{-2.5pt}
\begin{equation}
% \vspace{-3pt}
\begin{aligned}
    \text{MSE}(\gA(X,X_{true})) + \beta_1 D_{KL}(q_\phi(z | X, \gR)\parallel p_\psi(z|\gR)) + \beta_2 \frac{1}{|\gE^*|}\sum_{(i,j) \in\gE^*}||r_{ij} - r^{true}_{ij}||^2,
\label{eq:loss}
\end{aligned}
\end{equation}
where $\gA$ is the Kabsch alignment function~\citep{Kabsch}, $\gE^*$ are all the 1 and 2-hop edges in the molecular graph, with $r_{ij}$ corresponding to the distance between atoms $i$ and $j$. We note that both $\beta_1$ and $\beta_2$ play a crucial role in the optimization. $\beta_1$ has to be set low enough ($1e{-3}$) to allow the optimization to focus on the MSE when the differences between the model-based $X$ and the ground truth are very close, due to the RDKit distortion parameterization. 

For the QM9 experiments, $\beta_1$ is annealed starting from $1e{-6}$ to $1e{-1}$, increasing by a factor of 10 each epoch. $\beta_2$ controls the distance auxiliary loss and also had to be similarly annealed. We found that when $\beta_2 = 0$, \model still learned to improve upon the aligned MSE loss by 50\%, as compared to RDKit. Our error analysis showed that the resulting molecules either had extremely low distance error with high MSE, or vice-versa. Therefore, when the learning objective is unconstrained, our model learns to violate distance constraints by placing atoms in low-error but unphysical positions.

For QM9, by slowly annealing the distance loss, we allow our model to reach a metaphysical unstable transition state where distances are violated, but the aligned coordinate error is better. We then force the model to respect distance constraints. In the case of DRUGS, we found that this transition state was too difficult for the model to escape from, and we report the results using $\beta_2 = 0.5$ in \tref{tab:results_geom}. In Appendix \sref{section:appendix-more-drugs}, we further explore this idea and experiment with different annealing schedules for DRUGS. 
We note that as \model learns the torsion angles in an unsupervised manner because of the chosen CG strategy, we leave explicit angle optimization to future work.

\section{Coarse-graining}
\label{section:appendix-CG}
\vspace{-1ex}
We elaborate on the coarse-graining procedure introduced in \sref{sec:methods-algorithmic-details}. Following \citet{CGVAE}, we represent fine-grained (FG) molecular conformers as $x = \{x_i\}_{i=1}^n \in \R^{n \times 3}$. Similarly, the coarse-grained (CG) conformers are represented by $X = \{X_I\}_{I=1}^N \in \R^{N \times 3}$ where $N < n$.
Let $[n]$ and $[N]$ denote the set $\{1, 2, ..., n\}$ and $\{1, 2, ..., N\}$ respectively. 
The CG operation can be defined as an assignment $\map: [n] \rightarrow [N]$, which maps each FG atom $i$ in  $[n]$ to CG bead $I \in [N]$, i.e., bead $I$ is composed of the set of atoms $C_{I} = (k \in n \mid \map(k) = I )$. 
$X_I$ is initialized at the center of mass = $\frac{1}{|C_I|}\sum_{j \in C_I} x_j$.

We note that \model coarsens input molecules by first severing all torsion angles $\tau_{abcd}$, with $k$ torsion angles resulting in $k+1$ connected components or CG beads. This allows us, on average, to represent QM9 molecules with three beads and large drug molecules $(n>100)$ with 29 beads. We opted for a torsion angle-based strategy as it allows for unsupervised control over torsion angles, as well as the ability to rotate each subgraph independently. The CG strategy can be altered for various applications going forward.

\newpage
\section{Encoder Equations}
\vspace{-1ex}
\begin{figure*}[h]
\centering
    \includegraphics[width=\textwidth]{figures/encoder.png}
    \vspace{-10pt}
    \caption{\textbf{Encoder module message passing structure.} \textbf{(I)} Fine-grained graph with auxiliary 4\AA\ distance cut off. \textbf{(II)} Pooling graph with nodes for each atom and coarse-grained bead. Each group of nodes represents the formation of a CG bead. There is a single directional edge from each atom to its corresponding bead. \textbf{(III)} Coarse-grained graph with auxiliary 4\AA\ distance cut off, using the learned representation from the pooling graph.
    \model reduces the input from (I) to (III), drastically reducing the complexity of the problem.
    }
    \label{fig:encoder_graphs}
    \vspace{-5pt}
\end{figure*}
\subsection{Fine-grain Module}
\label{section:eqn-fine}
\vspace{-1ex}
% We expand on the fine-grained (FG) module described in \sref{section:hierarchy}. 
We describe the encoder, shown in~\fref{fig:schematic}(I). The model operates over SE(3)-invariant atom features $h \in R^{n\times D}$, and SE(3)-equivariant atomistic coordinates $x \in R^{n\times3}$. A single encoder layer is composed of three modules: fine-grained, pooling, and coarse-grained. Full equations for each module can be found in Appendix \sref{section:eqn-fine}, \sref{section:eqn-pool}, \sref{section:eqn-coarse}, respectively.

The fine-grained module is a graph-matching message-passing architecture. It differs from~\citet{equibind} by not having internal closed-form distance regularization and exclusively using unidirectional attention. It aims to effectively match the approximate conformer and ground truth by updating attention from the former to the latter.

The FG module is responsible for processing the FG atom coordinates and invariant features. More formally, the FG is defined as follows:
\begin{equation}
\begin{aligned}
\bm{m}_{j\to i} &= \phi^e(\bm{h}_i^{(t)},\bm{h}_j^{(t)},\norm{\bm{x}_i^{(t)}-\bm{x}_j^{(t)}}^2,\bm{f}_{j\to i}),\forall (I,J)\in \gE\cup\gE', \\
\bm{u}_{j'\to i} &= a_{j'\to i} \bm{W} \bm{h}_{j'}^{(t)},\forall i\in\gV,j'\in\gV', \\
\bm{m}_i &= \frac{1}{|\gN(i)|} \sum_{j\in\gN(i)}\bm{m}_{j\to i},\forall i\in\gV\cup\gV', \\
\bm{u}_{i} &=  \sum_{j'\in\gV'}\bm{u}_{j'\to i},\forall i \in\gV, \quad \text{and} \quad \bm{u}'_{i} =  0, \\
\bm{x}_i^{(t+1)} &= \eta_x \cdot \bm{x}_i^{(0)} + (1 - \eta_x) \cdot \bm{x}_i^{(t)} + \sum_{j\in\gN(i)} (\bm{x}_i^{(t)}-\bm{x}_j^{(t)})\phi^x(\bm{m}_{j\to i}), \\
% \bm{x}_i^{(t+1)} = \Psi\left(\bm{x}_i^{(t)} + \sum_{j\in\gN(i)}\frac{\bm{x}_i^{(t)}-\bm{x}_j^{(t)}}{\|\bm{x}_i^{(t)}-\bm{x}_j^{(t)} \|}\phi^x(\bm{m}_{j\to i})\right) \\
\bm{h}_i^{(t+1)} &= (1 - \eta_h) \cdot \bm{h}_i^{(t)} + \eta_h \cdot \phi^h(\bm{h}_i^{(t)},\bm{m}_i,\bm{u}_{i}, \bm{f}_i),\forall i\in\gV\cup\gV',
\end{aligned}
\label{fine-grain-eqn}
\end{equation}
where $f$ represents the original invariant node features $h^{t=0}$, $a_{j\to i}$ are SE(3)-invariant attention coefficients derived from $h$ embeddings, $\gN(i)$ are the graph neighbors of node $i$, and $W$ is a parameter matrix. $(\gV, \gE)$ and $(\gV', \gE')$ refer to the low-energy and RDKit approximation molecular graphs, respectively. The various $\phi$ functions are modeled using shallow MLPs, with  $\phi^x$ outputting a scalar and $\phi^e$ and $\phi^h$ returning a D-dimensional vector. $\eta_x$ and $\eta_h$ are weighted update parameters for the FG coordinates $x$ and invariant features $h$ respectively. We note that attention flows in a single direction from the RDKit approximation to the ground truth to prevent leakage in the parameterization of the learned prior distribution.  

\subsection{Pooling Module}
\label{section:eqn-pool}
% We expand on the pooling module described in \sref{section:hierarchy}. 
The pooling module takes in the updated representations ($h$ and $x$) of both the ground truth molecule and the RDKit reference from the FG module.
The pooling module is similar to the FG module, except it no longer uses attention and operates over a pooling graph.
Given a molecule with $n$ atoms and $N$ CG beads, the pooling graph consists of $n + N$ nodes. There is a single directional edge from all atoms to their respective beads. % (see~\fref{fig:novelty}(I)).
This allows message passing to propagate information through the predefined coarsening strategy.

The pooling module is responsible for learning the coordinates and invariant features of each coarse-grained bead by pooling FG information in a graph-matching framework. More formally, the pooling module is defined as follows:
\begin{equation}
\begin{aligned}
\bm{m}_{j\to I} &= \phi^e(\bm{H}_I^{(t)},\bm{h}_j^{(t)},\norm{\bm{X}_I^{(t)}-\bm{x}_j^{(t)}}^2,\bm{f}_{j\to I}),\forall (I,J)\in \gE\cup\gE', \\
\bm{m}_I &= \frac{1}{|\gN(I)|} \sum_{j\in\gN(I)}\bm{m}_{j\to I},\forall I\in\gV\cup\gV', \\
\bm{X}_I^{(t+1)} &= \eta_X \cdot \bm{X}_I^{(0)} + (1 - \eta_X) \cdot \bm{X}_I^{(t)} + \sum_{j\in\gN(I)} (\bm{X}_I^{(t)}-\bm{x}_j^{(t)})\phi^x(\bm{m}_{j\to I}), \\
% \bm{x}_i^{(t+1)} = \Psi\left(\bm{x}_i^{(t)} + \sum_{j\in\gN(i)}\frac{\bm{x}_i^{(t)}-\bm{x}_j^{(t)}}{\|\bm{x}_i^{(t)}-\bm{x}_j^{(t)} \|}\phi^x(\bm{m}_{j\to i})\right) \\
\bm{H}_I^{(t+1)} &= (1 - \eta_H) \cdot \bm{H}_I^{(t)} + \eta_H \cdot \phi^h(\bm{H}_I^{(t)},\bm{m}_I, \bm{f}_I),\forall I\in\gV\cup\gV',
\end{aligned}
\end{equation}
where capital letters refer to the CG representation of the pooling graph.
The pooling module mimics the FG module without attention on a pooling graph, as seen in \fref{fig:encoder_graphs}(II). The pooling graph contains a single node for each atom and CG bead, with a single edge from each FG atom to its corresponding bead. It is used to learn the appropriate representations of the CG information. As the pooling graph only contains edges from fine-to-coarse nodes, the fine-grain coordinates and features remain unchanged. The pooling graph at layer $t$ uses the invariant feature $H$ from the CG module of layer $t - 1$ to propagate information forward through the neural network. The main function of the pooling module is to act as a buffer between the FG and CG spaces. As a result, we found integrating the updated CG representation useful for building a better transition from FG to CG space.

\subsection{Coarse-grain Module}
\label{section:eqn-coarse}
The coarse-grained module uses the updated CG representations ($H\in R^{N\times D}$ and $X \in R^{N\times 3}$) from the pooling module to learn equivariant CG features ($Z$ and $\Tilde{Z} \in R^{N\times F \times 3}$) for the ground truth molecule and the RDKit reference. $F$ is fixed as a hyperparameter for latent space size. $N$ is allowed to be variable-length to handle molecules resulting from any coarsening procedure. 
The CG features are learned using a graph-matching point convolution~\citep{tensor-field-network} with similar unidirectional attention as the FG module. Prior to the main message-passing operations, the input features undergo equivariant mixing~\citep{3dlinker} to further distill geometric information into the learned CG representation.

The CG module is responsible for taking the pooled CG representation from the pooling module and learning a node-level equivariant latent representation.
We note that we use simple scalar and vector operations to mix equivariant and invariant features without relying on computationally expensive higher-order tensor products.
In the first step, invariant CG features $H$ and equivariant features $\bm{v}\in\mathbb{R}^{F\times 3}$ are transformed and mixed to construct new expressive intermediate features $H^{\prime}, H^{\prime\prime}, \bm{v}^{\prime}$ by,
\begin{subequations}
\begin{align}
H^{\prime}_I &= \phi_{1}(h^{(t)}_I, \norm{\text{VN-MLP}_{1}(\bm{v^{(t)}_I})})\in\mathbb{R}^{D} \label{eqn-a},\\ 
H^{\prime\prime}_I&=\phi_{2}(h^{(t)}_I, \norm{\text{VN-MLP}_{2}(\bm{v^{(t)}_I})})\in\mathbb{R}^{F}\label{eqn-b},\\
\bm{v}^{\prime}_I &=\text{diag}\{\phi_{3}(H^{(t)}_I)\}\cdot\text{VN-MLP}_{3}(\bm{v^{(t)}_I})\in \mathbb{R}^{F\times 3}.\label{eqn-c}
\end{align}
\label{mixed features}
\end{subequations}
Next, a point convolution \citep{tensor-field-network, painn, 3dlinker} is applied to linearly transform the mixed features $H^{\prime},H^{\prime\prime},\bm{v}^{\prime}$ into messages:
\begin{subequations}
\begin{align}
\bm{m^H}_{I\leftarrow J} &= \text{Ker}_{1}(\norm{\bm{r}_{I,J}})\odot H^{\prime}_J\label{eqn-d}, \\
\bm{m^v}_{I\leftarrow J}&=  \text{diag}\left\{\text{Ker}_{2}(\norm{\bm{r}_{I,J}})\right\} \cdot \bm{v}^{\prime}_J  +\big(\text{Ker}_{3}(\norm{\bm{r}_{I,J}})\odot H^{\prime\prime}_{J}\big)\cdot\bm{r}^{\top}_{I,J}\label{eqn-e}, \\
\bm{u}_{J'\to I} &= a_{J'\to I} \bm{W} \bm{H}_{J'}^{(t)},\forall I\in\gV,J'\in\gV', \\
\bm{u}_{I} &=  \sum_{J'\in\gV'}\bm{u}_{J'\to I},\forall I \in\gV, \quad \text{and} \quad \bm{u}'_{I} =  0, \\
\bm{H}^{t+1}_{I}& = (1 - \eta_H) \cdot H^{\ell}_I+ \eta_H \cdot\text{MLP}(H^{\ell}_I, \sum_{J\in N(I)}m^H_{I\leftarrow J}, u_{I}),\forall I\in\gV\cup\gV' \label{agg-h},\\
\bm{v}^{t+1}_{I} &= (1 - \eta_v) \cdot v^{\ell}_I + \eta_v \cdot \text{VN-MLP}_{4}(\bm{v}^{\ell}_I, \sum_{J\in N(I)}\bm{m^v}_{I\leftarrow J}),\forall I\in\gV\cup\gV'\label{agg-v},
\end{align}
\label{message function}
\end{subequations}
where each Ker refers to a learned RBF kernel, $r_{IJ}$ is the difference between $X_I$ and $X_J$, and $a_{J\to I}$ are SE(3)-invariant attention coefficients derived from the learned invariant features $H$. $\eta_H$ and $\eta_v$ control the mixing of the learned invariant and equivariant representations.

We note that for $t > 0$, the $H_I$ from the CG module are used in the next layer's pooling module, creating a cyclic dependency to learn an information-rich CG representation. This is shown by the dashed lines in~\fref{fig:schematic}(I). The cyclic flow of information grounds the learned CG representation to the innate FG structure. All equivariant CG features $v$ are initialized as zero and are slowly built up through each message passing layer. As point convolutions and VN operations are strictly SO(3)-equivariant, we subtract the molecule’s centroid from the atomic coordinates prior to encoding, making it effectively SE(3)-equivariant.

The modules in each encoder layer communicate with the respective module of the previous layer. This hierarchical message-passing scheme results in an informative and geometrically grounded final CG latent representation. We note that the pooling module of layer $\ell$ uses the updated invariant features $H$ from the CG module of layer $\ell - 1$, as shown by the dashed lines in~\fref{fig:schematic}(I).

\section{Decoder Architecture}
\label{section:decoder}
\vspace{-1ex}

We sample from the learned posterior (training) and learned prior (inference) to get $Z = \mu + \eps \sigma$, where $\eps$ is noise sampled from a standard Gaussian distribution as the input to the decoder. 
We note the role of the decoder is two-fold. The first is to convert the latent coarsened representation back into FG space through a process we call channel selection. The second is to refine the fine-grain representation autoregressively to generate the final low-energy coordinates.

\vspace{-1ex}
\paragraph{Channel Selection.} To explicitly handle all choices of coarse-graining techniques, our model performs variable-length backmapping.
% To improve conformer generation using coarse-graining, we need to ensure that our model can handle variable-length backmapping. 
This aspect is crucial because every molecule can be coarsened into a different number of beads, and there is no explicit limit to the number of atoms a single bead can represent. Unlike CGVAE~\citep{CGVAE}, which requires training a separate model for each choice in granularity $N$, \model is capable of reconstructing FG coordinates from any $N$ (illustrated in \fref{fig:schematic}(III)).

CGVAE defines the process of channel selection as selecting the top $k$ latent channels, where $k$ is the number of atoms in a CG bead of interest. Instead of discarding all learned information in the remaining $F-k$ channels in the latent representation, we use a novel aggregated attention mechanism. This mechanism learns the optimal mixing of channels to reconstruct the FG coordinates and is illustrated in~\fref{fig:attention}.
The attention operation allows us to actively query our latent representation for the number of atoms we need, and draw upon similarities to the learned RDKit approximation that has been distilled into the latent space through the encoding process.
Channel selection translates the CG latent tensor $Z \in R^{N\times F \times 3}$ into FG coordinates $x_{cs} \in R^{n\times 3}$.

\vspace{-1ex}
\paragraph{Coordinate Refinement.} 
Once channel selection is complete, we have effectively translated the variable-length CG representation back into the desired FG form. From here, $x_{cs}$ is grouped into its corresponding CG beads but left in FG coordinates to do a bead-wise autoregressive generation of final low-energy coordinates (\fref{fig:schematic}(IV)). 
As there is no intrinsic ordering of subgraphs, we use a breadth-first search that prioritizes larger subgraphs with large out-degrees. In other words, we generate a linear order that focuses on the largest, most connected subgraphs and works outward. We believe that by focusing on the most central component first, which occupies the most 3D volume, we can reduce the propagation of error that is typically observed in autoregressive approaches.
We stress that by coarse-graining by torsion angle connectivity, our model learns the optimal torsion angles in an unsupervised manner, as the conditional input to the decoder is not aligned. 
\model ensures each next generated subgraph is rotated properly to achieve a low coordinate and distance error.

\vspace{-1ex}
\paragraph{Learning the Optimal Distortion.}The decoder architecture is similar to the EGNN-based FG layer in the encoder. However, it differs in two important ways. First, we mix the conditional coordinates with the invariant atom features using a similar procedure as in the CG layer instead of typical graph matching. Second, we learn to predict the difference between the RDKit reference and ground truth conformations. This provides an upper error bound and enables us to leverage easy-to-obtain approximations more effectively.

More formally, a single decoder layer is defined as follows:
\begin{subequations}
\begin{align}
\bm{\mu}^{(t)} &= \frac{1}{|\gV_{prev}|} \sum_{k\in \gV_{prev}}x_k, \\
\bm{\Tilde{h}}_{i} &= \phi^m(\bm{h}_i^{(t)},\bm{x}_i^{(t)},\bm{\mu}^{(t)},\norm{\bm{x}_i^{(t)}-\bm{\mu}^{(t)}}^2),\forall i \in \gV_{cur}, \\
\bm{m}_{j\to i} &= \phi^e(\bm{\Tilde{h}}_i^{(t)},\bm{\Tilde{h}}_j^{(t)},\norm{\bm{x}_i^{(t)}-\bm{x}_j^{(t)}}^2,\norm{\bm{x}_i^{(t)}-\bm{x}_{ref, j}^{(t)}}^2, \norm{\bm{x}_i^{(t)}-\bm{x}_{ref, i}^{(t)}}^2),\forall (i,j)\in \gE_{cur}, \\
\bm{m}_i &= \frac{1}{|\gN(i)|} \sum_{j\in\gN(i)}\bm{m}_{j\to i},\forall i\in\gV_{cur}, \\
\bm{u}_{j'\to i} &= a_{j'\to i} \bm{W} \bm{h}_{j'}^{(t)},\forall i\in\gV_{cur},j'\in\gV_{prev}, \\
\bm{u}_{i} &=  \sum_{j'\in\gV_{prev}}\bm{u}_{j'\to i},\forall i \in\gV_{cur},\\
\bm{x}_i^{(t+1)} &= \bm{x}_{ref,i}^{(t)} + \sum_{j\in\gN(i)} (\bm{x}_i^{(t)}-\bm{x}_j^{(t)})\phi^x(\bm{m}_{j\to i}),\forall i \in\gV_{cur}, \\
\bm{h}_i^{(t+1)} &= (1 - \beta) \cdot \bm{h}_i^{(t)} + \beta \cdot \phi^h(\bm{\Tilde{h}}_i^{(t)},\bm{m}_i,\bm{u}_{i}, \bm{f}_i),\forall i\in\gV_{cur},
\end{align}
\label{eqn:decoder}
\end{subequations}
where ($\gV_{cur}$, $\gE_{cur}$) and ($\gV_{prev}$, $\gE_{prev}$) refer to the subgraph currently being generated and the set of all previously generated subgraphs, \ie, the current state of the molecule. $\phi^m, \phi^e, \phi^x, \text{and } \phi^h$ refer to separate shallow MLPs for the feature mixing, edge message calculation, coordinate update, and invariant feature update, respectively. \eref{eqn:decoder}(a-b) creates a mixed feature for each atom consisting of the current FG invariant feature and 3D position vectors ($h$ and $x$), and the previous centroid $\mu$ and respective centroid distances. \eref{eqn:decoder}(c-d) defines the message passing operation that uses the aforementioned mixed features $\Tilde{h}$ and a series of important distances between the model-based conformer and RDKit reference.% (see~ \fref{fig:novelty}(II) for a helpful visualization). 
\eref{eqn:decoder}(e-f) apply the same unidirectional attention updates seen in the encoder architecture. \eref{eqn:decoder}(g-h) update the position and feature vector for each atom using the above messages and attention coefficients, with $f$ representing the original invariant node features $h^{\ell=0}$ and $\beta$ a weighted update parameter. We emphasize that~\eref{eqn:decoder}(g) formulates the overall objective as learning the optimal distortion of the RDKit reference to achieve the low-energy position \ie,\ $x^* = x_{ref} + \Delta x$.
The CG autoregressive strategy allows \model to handle extremely large molecules efficiently, as the max number of time steps is equal to the max number of CG beads. \model is trained using teacher forcing~\citep{teacher-forcing}, which enables an explicit mixing of low-energy coordinates with the current FG positions from channel selection \eref{eqn:decoder}(a-b).
\vspace{-2ex}

\section{Model Configuration}
\label{section:appendix-breakdown}
\vspace{-1ex}
\paragraph{Model.} We present the model configuration that was used to generate the results in \sref{sec:rmsd} - \sref{sec:rigid}.
Overall, the default model has 1.9M parameters: 1.6M for the encoder and 300K for the decoder. We note that as \model uses graph matching, half the encoder parameters are used for each of the two inputs representing the same molecule in different spatial orientations. For both the encoder and decoder, we use five message-passing layers, a learning rate of $1e{-3}$ with an 80\% step reduction after each epoch, and a latent space channel dimension ($F$) of 32. All other architectural parameters, such as feature mixing ratios or nonlinearities, were set following similar architectures~\citep{3dlinker, vn-mlp, equibind}. We present further ablations in Appendix \sref{section:appendix-more-drugs}. We note that the ability to share weights between the inputs as well as between each layer in the encoder is left as a hyperparameter. This could allow the encoder to see a 2x or 5x reduction in model size, respectively.

\paragraph{Compute.}
The QM9 model was trained and validated for five epochs in 15 hours using a single 40GB A100 GPU. We used a batch size of 600, where a single input refers to two graphs: the ground truth and RDKit approximate conformer. The DRUGs model was trained and validated for five epochs in 50 hours using distributed data-parallel (DDP) with 4 40GB A100 GPUs with a batch size of 300 on each GPU. For DRUGs, the GPU utilization was, on average, 66\% as few batches contain very large molecules. In the future, lower run times can be achieved if the large molecules are more intuitively spaced out in each batch. 

We note DDP has a negative effect on overall model benchmark performance due to the gradient synchronization but was used due to compute constraints. Without DDP, we expect the training time to take around 7 days, which is on par with Torsional Diffusion (4-11 days). 
We demonstrated that \model achieves as good or better results than prior methods with less data and time, and these results can be further optimized in future work. We provide evidence of the negative effects of DDP in Appendix \sref{section:appendix-more-drugs}.

\paragraph{Optimal Transport reduces compute requirements.} The optimal transport (OT) models were trained on 2 epochs on a single A6000 GPU for 8 and 15 hours total for QM9 and DRUGS, respectively. For OT details, see Appendix \eref{eq:ot_loss}. Here, both models use the first 5 ground truth conformers.
In real-world applications like polymer design, the availability of data is frequently limited and accompanied by a scarcity of conformers for each molecule. The current datasets, QM9 and DRUGS, do not mimic this setting very well. For example, on average, QM9 has 15 conformers per molecule, and DRUGS has 104 per molecule---both datasets have significantly more conformers than in an experimental drug design setting. Given this, rather than training on the first 30 conformers as done in Torsional Diffusion, we train on the first five (typically those with the largest Boltzmann weight) for QM9 and DRUGS, respectively. 
% This corresponds to 41\% of TD's training set or 23\% of the overall dataset for QM9, and 7.3\% of TD's training set or 1.7\% of DRUGS overall.
% Even with less data, we outperform models trained on significantly larger training sets.

\section{RDKit Approximate Conformers}
\label{section:mmff}
\paragraph{Generating Approximate Conformers.} For \modelns's initial conditional approximations, we only use RDKit + MMFF when it can converge (\textasciitilde{90\%} and \textasciitilde{40\%} convergence for QM9 and DRUGS, respectively). We emphasize that RDKit only throws an error when MMFF is not possible but often returns structures with a non-zero return code, which signifies incomplete and potentially inaccurate optimizations. Therefore, in generating the RDKit structures for training and evaluation, we filter for MMFF converged structures. We default to the base EKTDG-produced structures when either the optimization cannot converge, or MMFF does not yield enough unique conformers. \model ultimately offers a solution that can effectively learn from traditional cheminformatics methods. This aspect of MMFF convergence has not been discussed in prior ML for MCG methods, and we leave it to future cheminformatics research to learn the causes and implications of incomplete optimizations.

\paragraph{Eliminating distribution shift with explicit conditioning.}
Both \model and TD optimize $p(X|\gR)$ but utilize the RDKit approximations $\gR$ in different ways. TD learns to update the torsion angles of $\gR$, while \model leverages CG information to inform geometric updates (coordinates, distances, and torsion angles) to translate $\gR$ to $X$.
Unlike TD, which uses a preprocessing optimization procedure to generate substitute ground truth conformers that mimic $p(\gR)$, 
\model directly learns from both $X$ and $\gR$ through its hierarchical graph matching procedure. This directly addresses the distributional shift problem.
We hypothesize that this, along with our angle-based CG strategy, leads to our observed improvements.
Overall, \model provides a comprehensive framework for accurate conformer generation that can be directly applied for downstream tasks such as oracle-based protein docking. 

\section{GEOM Benchmark Discussion}
\label{section:geom-xtb}
\paragraph{xTB energy and property prediction.}
% To directly assess the quality of generated samples, we estimate the molecular ensemble properties~\citep{GEOM} over a set of generated conformers. 
% Specifically, similar to TD and GD, we sample min(2K, 32) conformers of 100 GEOM-DRUGS molecules, where K is the number of ground truth conformers. 
% All structures are relaxed with GFN2-xTB~\citep{xtb} with the BGFS optimizer before predicting Energy $E$, HOMO-LUMO Gap $\Delta \epsilon$, and minimum Energy $E_{min}$  in kcal/mol and dipole moment $\mu$ in debye. We report the mean absolute error of generated properties against the ground truth ensemble properties.
We note the issues surrounding the RMSD metrics have always existed, and prior MCG methods have introduced energy-based benchmarks that we describe and report in~\tref{tab:xtb_energy}. We note these energies are calculated with xTB, and thus are not very accurate compared to density functional theory (DFT), as it is limited by the level of theory used to produce the energies further discussed in \cite{GEOM}.
Therefore, since current benchmarks mainly focus on gauging the effectiveness of the machine learning objective and less on the chemical feasibility and downstream use of the generated conformers, we use oracle-based protein docking-based to evaluate conformer quality on downstream tasks.
These evaluations are highly informative, as molecular docking is a crucial step in the drug discovery process, as it helps researchers identify potential drug candidates and understand how they interact with their target proteins.
The combination of RMSD, xTB energy, and downstream docking tasks presents a more comprehensive evaluation of generated conformers.
\section{QM9 Experimental Details}
\label{section:appendix-qm9}
\vspace{-1ex}
% The QM9 test set was normalized to ensure all methods were evaluated on the exact same molecules due to a small number of preprocessing failures. These issues were tied to RDKit failures and the use of different definitions of torsion angles. The normalization enforces an equal success rate and results in only an overall dataset reduction of 3\%. 
% We note for QM9 specifically, RDKit without MMFF achieves better recall than with MMFF but worse precision. We leave it to future work to explore the implications of this and how to best take advantage of built-in optimization methods.

Both \model and \modelns-OT were trained on 5 conformers per ground truth molecule, compared to Torsional Diffusion's 30. We hypothesize that since \model uses a one-to-one loss function, we are able to maintain high recall, whereas the OT model finds an optimal matching that focuses on precision. By adding more ground truth conformers, we hypothesize our model can better cover the true conformer space, improving recall, as the OT setting would not be as biased toward precision. 

\begin{table}[H]
\centering
\caption{Quality of generated conformer ensembles for the GEOM-QM9 test set ($\delta=0.5$\AA) in terms of Coverage (\%) and Average RMSD (\AA). Torsional Diffusion (TD) was benchmarked using its evaluation code and available generated molecules, per their public instructions. Note that \model (5 epochs) was restricted to using 41\% of the data used by TD (250 epochs) to exemplify a low-compute and data-constrained setting. OMEGA results were taken from \cite{torsional-diffusion} (we were unable to run the coverage normalization).}
\label{tab:results_qm9_supp}
\begin{tabular}{l|cccc|cccc} \toprule
               & \multicolumn{4}{c|}{Recall} & \multicolumn{4}{c}{Precision} \\
                  & \multicolumn{2}{c}{Coverage $\uparrow$} & \multicolumn{2}{c|}{AR $\downarrow$} & \multicolumn{2}{c}{Coverage $\uparrow$} & \multicolumn{2}{c}{AR $\downarrow$} \\
Method  &  Mean & Med & Mean & Med & Mean & Med & Mean & Med \\ \midrule
\multicolumn{1}{l|}{OMEGA} &   85.5 &100.0 &0.177 &0.126 & 82.9 & 100.0 & 0.224 &0.186\\
\multicolumn{1}{l|}{RDKit + MMFF} &   75.2 & 100.0 & 0.219 & 0.173 & 82.1 & 100.0 & 0.157 & 0.119\\
% TD no filter  88.6 100 0.177 0.146 84.6 100 0.219 0.194
\multicolumn{1}{l|}{GeoMol} & 79.4 & 100.0 & 0.219 & 0.191 & 75.9 & 100.0 & 0.262  & 0.233\\
\multicolumn{1}{l|}{Torsional Diffusion} & 82.2 & 100.0 & 0.179 & 0.148 & 78.4 & 100.0 & 0.222  & 0.197\\
% \model (min) & 81.6 & \textbf{100.0} & 0.285 & 0.249 & 88.2 & \textbf{100.0} & 0.236 & 0.194 \\ no normalization filtering
\multicolumn{1}{l|}{\model} &  76.9 & 100.0 & 0.246 & 0.211 & 80.2 & 100.0 & 0.227 & 0.186\\
\multicolumn{1}{l|}{\modelns-OT} &  56.1 & 50.0 & 0.361 & 0.345 & 80.2 & 100.0 & 0.149 & 0.108\\
\bottomrule
\end{tabular}
\end{table}

\section{DRUGS Extended Benchmarks}
\label{section:appendix-more-drugs}
\vspace{-1ex}
% \danny{fix this section commenting out for now}
% \begin{table}[H]
% \begin{center}
% \caption{Quality of generated conformer ensembles for the GEOM-DRUGS test set ($\delta=0.75$\AA) in terms of Coverage (\%) and Average RMSD (\AA). \model was trained for 5 epochs unless specified. \model$\dagger$ refers to training for 0.2 epochs on a single GPU with 5 conformers per training molecule and no DDP. We provide GeoDiff results using the same generated molecules as in \citet{torsional-diffusion}. We provide GeoMol results using its published weights and code \citep{GeoMol}. 
% % In terms of maintaining equivariance, all \model runs use scalar and vector operations, which is analogous to $\ell = 1$. 
% % In terms of maintaining equivariance, both \model and TD ($\ell=1$) use scalars and vectors to create their respective equivariant bases.
% Both \model and TD ($\ell$ = 1) use scalar and vector operations to maintain equivariance. 
% TD ($\ell=2$) additionally uses higher-order tensor products (see Appendix \sref{section:appendix-limits} for more details).}
% \label{tab:results_drugs_app}
% \begin{tabular}{lc|cccc|cccc} \toprule
%                & \multicolumn{1}{l|}{} & \multicolumn{4}{c|}{Recall} & \multicolumn{4}{c}{Precision}  \\
%                  & \multicolumn{1}{l|}{} & \multicolumn{2}{c}{Coverage $\uparrow$} & \multicolumn{2}{c|}{AR $\downarrow$} & \multicolumn{2}{c}{Coverage $\uparrow$} & \multicolumn{2}{c}{AR $\downarrow$} \\
% Method  & Metric & Mean & Med & Mean & Med & Mean & Med & Mean & Med \\ \midrule
% \multicolumn{1}{l|}{RDKit} & min  & 39.7 & 28.9 & 1.052 & 0.989 & 41.9 & 32.9 & 0.988 & 0.878 \\
% \multicolumn{1}{l|}{RDKit + MMFF} & min  & 49.0 & 46.2 & 0.896 & 0.825 & 55.3 & 55.1 & 0.840 & \textbf{0.715} \\
% \multicolumn{1}{l|}{GeoMol} & min  & 45.3 & 42.8 & 0.866 & 0.826 & 43.4 & 42.8 & 0.919 & 0.842 \\
% % \multicolumn{1}{l|}{GeoDiff} & min  & 32.0 & 18.0 & 0.805 & 0.786 & 20.4 & 6.89 & 1.067 & 1.016 \\
% \multicolumn{1}{l|}{GeoDiff} & min  & 43.1 & 39.3 & 0.828 & 0.806 & 25.5 & 15.2 & 1.131 & 1.083 \\
% \multicolumn{1}{l|}{Torsional Diffusion ($\ell = 1$)} & min & 71.8 & \textbf{80.8} & 0.591 & 0.568 & 52.5 & 52.7 & 0.804 & 0.758 \\
% \multicolumn{1}{l|}{Torsional Diffusion ($\ell = 2$)} & min & \textbf{73.5} & 80.6 & \textbf{0.574} & \textbf{0.552} & \textbf{55.9} & \textbf{57.6} & \textbf{0.770} & 0.720 \\
% \multicolumn{1}{l|}{\model (2 conf)} & min & 45.7 & 42.4 & 0.888 & 0.848 & 47.0 & 41.8 & 0.915 & 0.823 \\
% \multicolumn{1}{l|}{\model (5 conf)} & min & 43.2 & 39.2 & 0.907 & 0.868 & 42.9 & 42.7 & 0.970 & 0.909 \\
% \multicolumn{1}{l|}{\model$\dagger$} & min & 16.4 & 0.00 & 1.147 & 1.123 & 14.2 & 0.00 & 1.252 & 1.212 \\
% \midrule
% \multicolumn{1}{l|}{RDKit} & mean  & 8.07 & 0.00 & 1.960 & 1.912 & 7.18 & 0.00 & 1.960 & 1.912 \\
% \multicolumn{1}{l|}{RDKit + MMFF} & mean & \textbf{8.87} & 0.00& 1.903 & 1.887 & \textbf{8.34} & 0.00& 1.903 & 1.887 \\
% \multicolumn{1}{l|}{GeoMol} & mean  & 7.27 & 0.00 & 1.847 & 1.892 & 6.89 & 0.00 & 1.847 & 1.892 \\
% % \multicolumn{1}{l|}{GeoDiff} & mean  & 3.46 & 0.00 & 1.906 & 1.968 & 3.65 & 0.00 & 1.906 & 1.968 \\
% \multicolumn{1}{l|}{GeoDiff} & mean  & 3.64 & 0.00 & 1.985 & 2.033 & 4.02 & 0.00 & 1.985 & 2.033 \\
% \multicolumn{1}{l|}{Torsional Diffusion ($\ell = 1$)} & mean & 7.16 & 0.00 & 1.804 & 1.873 & 7.48 & 0.00 & 1.804 & 1.873 \\
% \multicolumn{1}{l|}{Torsional Diffusion ($\ell = 2$)} & mean & 7.19 & 0.00 & 1.793 & 1.864 & 7.90 & 0.00 & 1.793 & 1.864\\
% \multicolumn{1}{l|}{\model (2 conf)} & mean & 8.51 & 0.00 & 1.791 & 1.825 & 7.83 & 0.00& 1.791 & 1.825 \\
% \multicolumn{1}{l|}{\model (5 conf)} & mean & 8.29 & 0.00 & 1.778 & 1.815 & 7.81 & 0.00 & 1.778 & 1.815 \\
% \multicolumn{1}{l|}{\model$\dagger$} & mean & 5.00 & 0.00 & \textbf{1.777} & \textbf{1.786} & 4.81 & 0.00 & \textbf{1.777} & \textbf{1.786} \\
% \midrule
% \multicolumn{1}{l|}{RDKit} & max  & 3.62 & 0.00 & 2.603 & 2.580 & 3.56 & 0.00 & 2.777 & 2.804 \\
% \multicolumn{1}{l|}{RDKit + MMFF} & max & \textbf{4.16} & 0.00& 2.589 & 2.551 & 3.95 & 0.00& 2.734 & 2.774 \\
% % \multicolumn{1}{l|}{GeoDiff} & max  & 1.60 & 0.00 & 2.837 & 2.938 & 2.11 & 0.00 & 2.615 & 2.733 \\
% \multicolumn{1}{l|}{GeoMol} & max  & 3.23 & 0.00 & 2.739 & 2.804 & 3.73 & 0.00 & 2.633 & 2.727 \\
% \multicolumn{1}{l|}{GeoDiff} & max  & 1.66 & 0.00 & 2.946 & 3.062 & 2.28 & 0.00 & 2.711 & 2.806 \\
% \multicolumn{1}{l|}{Torsional Diffusion ($\ell = 1$)} & max & 2.97 & 0.00 & 2.840 & 2.947 & 3.81 & 0.00 & 2.617 & 2.717 \\
% \multicolumn{1}{l|}{Torsional Diffusion ($\ell = 2$)} & max & 3.44 & 0.00 & 2.813 & 2.930 & \textbf{3.96} & 0.00 & 2.615 & 2.726 \\
% \multicolumn{1}{l|}{\model (2 conf)} & max & 3.97 & 0.00 & 2.481 & 2.505 & 3.85 & 0.00 & 2.521 & 2.602 \\
% \multicolumn{1}{l|}{\model (5 conf)} & max & 4.07 & 0.00 & 2.471 & 2.495 & 3.87 & 0.00 & 2.474 & 2.578 \\
% \multicolumn{1}{l|}{\model$\dagger$} & max & 2.50 & 0.00 & \textbf{2.403} & \textbf{2.411} & 2.80 & 0.00 & \textbf{2.347} & \textbf{2.404} \\
% \bottomrule
% \end{tabular}
% \end{center}
% \end{table}
% We report ablation results for the DRUGS dataset in \tref{tab:results_drugs_app}\footnote{GeoMol failed on 0.1\% more molecules than \model and Torsional Diffusion.}.
% We provide RDKit with and without MMFF, GeoMol, GeoDiff, Torsional Diffusion (TD) ($\ell=1$ and $\ell = 2$), and \model trained on 2 and 5 conformers per training molecule, as well as without DDP. We see that for the chosen threshold $\delta = 0.75$, the addition of MMFF to RDKit results in minor improvements across all metrics. Further coverage results for all tested thresholds can be seen in \fref{fig:one}-\fref{fig:three}. 
\paragraph{Evaluation Details.} All models in \tref{tab:results_geom} were benchmarked with Torsional Diffusion's (TD) evaluation code and retrained if generated molecules were not public (using their public instructions). We note that TD uses higher-order tensor products to maintain equivariance ($\ell = 2$). In contrast, GeoMol, GeoDiff, and \model use scalar-vector operations that are theoretically analogous to $\ell = 1$. \modelns-OT uses an optimal transport (OT) loss with the same decoder architecture as in \fref{fig:schematic}, but is no longer autoregressive. GeoDiff's code would not load, so we were able to evaluate the GeoDiff generated DRUGS molecules from the Torsional Diffusion authors' evaluation on the same test set.

\begin{table}[H]
\centering
\caption{DRUGS-Precision equivariance ablations. OMEGA~\citep{omega} results were taken from~\cite{torsional-diffusion}. All others were re-benchmarked using Torsional Diffusion's code with an error normalized Coverage score to prevent the masking out of method failures. This enforces that each method is fairly evaluated on the entire test set as now Coverage truly represents the percentage of the test set that meets the threshold criteria. OMEGA requires a commercial license, so we were unable to test the results ourselves, thus taking results from TD. As a non-ML method, we also assume OMEGA has no failures as each molecule in the test set is valid, which could artificially inflate the observed coverage scores.} %We also include the results for RDKit+MMFF, but force all torsion angles to be equal to that of the ground truth conformer as a theoretical thought experiment to better understand the geometric qualities that impact the final results.
\begin{tabular}{lc|cccc} \toprule
                % & \multicolumn{4}{c}{DRUGS-Precision}  \\
                 & \multicolumn{1}{l|}{} & \multicolumn{2}{c}{Coverage $\uparrow$} & \multicolumn{2}{c}{AMR $\downarrow$} \\
Method  &  & Mean & Med & Mean & Med  \\ \midrule
% \multicolumn{2}{l|}{RDKit + MMFF} &    \textbf{82.1} & \textbf{100.0} & \textbf{0.157} & \textbf{0.119} & 55.3 & 55.1 & 0.840 & \textbf{0.715} \\
% RDKit would need to be rerun its removal boosts performance
% \multicolumn{2}{l|}{RDKit+MMFF+Perfect Angles}  & 71.3 & 89.3 & 0.515 & 0.428 \\
% \midrule
\multicolumn{2}{l|}{RDKit}  & 37.9 & 29.9 & 0.988 & 0.878 \\
\multicolumn{2}{l|}{RDKit + MMFF} & 52.3 & 52.1 & 0.840 & 0.715\\
\multicolumn{2}{l|}{OMEGA} &    53.4 & 54.6 & 0.841  & 0.762 \\
% \midrule
\multicolumn{2}{l|}{GeoDiff} &    23.7 & 13.0 & 1.131 & 1.083 \\
\multicolumn{2}{l|}{GeoMol} &    40.5 & 33.5 & 0.919 & 0.842 \\
\multicolumn{2}{l|}{Torsional Diffusion ($\ell = 1$)} &  48.9 & 50.0 & 0.804 & 0.758 \\
\multicolumn{2}{l|}{Torsional Diffusion ($\ell = 2$)} &  52.1 & 53.7 & 0.770 & 0.720 \\
\multicolumn{2}{l|}{\model} & 43.8 & 35.5 & 0.914 & 0.829 \\
\multicolumn{2}{l|}{\modelns-OT} & 52.0 & 52.1 & 0.836 & 0.694 \\
\bottomrule
\end{tabular}
\label{tab:results_geom_supp}
\vspace{-3ex}
\end{table}

We copy the results from \tref{tab:results_geom} and provide additional results, including TD for rotation order $\ell = 1$, OMEGA~\citep{omega}, and RDKit. This allows for a closer comparison to the scalar and vector operations that \model employs to maintain equivariance. Using a lower rotation order results in slightly worse results in nearly all categories. We further discuss the implications of the choice in equivariant representation in Appendix \sref{section:appendix-limits}.

\paragraph{Optimal Transport.} In practice, our model generates a set of conformers, $\{\mathcal{C}_k \}_{k \in [1..K]}$, that needs to match a variable-length set of low-energy ground truth conformers,  $\{\mathcal{C}^*_l\}_{l \in [1..L]}$. In our case, the number L of true conformers, or the matching between generated and true conformers is not known upfront.
For these reasons, we introduce an optimal transport-based, minimization-only, loss function~\citep{GeoMol}:
\begin{equation}
\begin{aligned}
    \mathcal{L}_{OT} &=  \min_{\T\in\mathcal{Q}_{K,L}} \sum_{k,l}T_{kl} \mathcal{L}(\mathcal{C}_k, \mathcal{C}^*_l), \\ 
    \mathcal{L}(\mathcal{C}_k, \mathcal{C}^*_l) &= \text{MSE}(\mathcal{C}_k, \mathcal{C}^*_l) + \text{distance error}(\mathcal{C}_k, \mathcal{C}^*_l),
    \label{eq:ot_loss}
\end{aligned}
% \vspace{-2ex}
\end{equation}
where $\T$ is the \textbf{\textit{transport plan}} satisfying $\mathcal{Q}_{K,L} = \{\T \in \R_+^{K\times L} : \T \mathbf{1}_L=\frac{1}{K}\mathbf{1}_K, \T^T \mathbf{1}_K=\frac{1}{L}\mathbf{1}_L \}$.
The minimization w.r.t. T is computed quickly using the Earth Mover Distance and the POT library~\citep{OT}. As the OT loss focuses more on finding the optimal mapping from generated conformers to ground truth reference, we removed the autoregressive decoding path of \model and replaced it with a single pass with the same decoder architecture. The underlying loss function, which is tasked to minimize MSE coordinate error, and interatomic distance error is the same in both (\eref{eq:loss}), the autoregressive (AR) and non-AR OT-based loss functions. The OT version additionally finds the optimal mapping between the generated and ground truth structures, which better aligns with the AMR and Coverage benchmarks.

% We provide several different \model runs with different hyperparameters.
% We saw that with constant training iterations, by increasing the dataset from 2 to 5 conformers per molecule, we get a low RMSD (AR) with slight decreases in the min results and mean coverage. 
% We hypothesize that by increasing the data by 2.5X, more epochs of training are necessary to increase the coverage, but we show improved error and generalization with more data. 
% We note all \model models were trained with the exact same hyperparameters and faced the observed DDP issues as discussed in Appendix \sref{section:appendix-breakdown}. 
% The no DDP result with only 0.2 of an epoch of training with 5 conformers shows how gradient synchronization of distributed training really hurt performance. We see that for mean, and max AR, the no DDP model achieves the best results, albeit with low coverage. 
% Training longer with more data will likely improve the coverage values. 
% This example yields an interesting perspective on the \model learning framework as a whole. 
% These results demonstrate that the absence of gradient synchronization helps greatly. In this case, the lowest error is achieved with very little data and training iterations

% % TODO explore differences between kabsch and getbestRMS
% %Furthermore, we found that by replacing RDKit's GetBestRMS function, which is used in the benchmarking procedure~\citep{torsional-diffusion} to calculate the lowest RMSD for all possible alignments between the reference and generated conformer\footnote{Hydrogen atoms are removed before applying the GetBestRMS function} with the Kabsch alignment (used in Appendix \sref{section:appendix-loss-func}), the AR error for the min, mean, and max metrics increases  by 0.6-1.0. This is significant as \modelns's loss function is not capturing what we intend the model to be used for. We leave to future work to explore this difference and bridge the gap between RDKit's expensive alignment function and differentiable loss functions.#Note BestRMS uses K^2 SVD operations

\paragraph{Hyperparameter Ablations.} We experimented with increasing the latent channels ($F$) from 32 to 64 and 128, and introducing a step-wise distance loss and KL regularization annealing schedule, as done in the QM9 experiments. Both these experiments resulted in slightly worse performance when limited to 2 conformers per training molecule. We hypothesize that due to the DRUGs molecules being much larger than those in QM9, more training may be necessary, and a more sensitive annealing schedule may be required.
% We note that the number of latent channels directly impacts the dimensionality of the aggregated attention operation (\fref{fig:attention}). A larger latent space may have more information, but may make the variable length coarse-to-fine translation more difficult.


\begin{figure}[htb]%[b]
  \centering
  \begin{subfigure}{0.45\textwidth}
    \centering
    \includegraphics[width=\textwidth]{figures/rec_cov3.png}
    \caption{Recall Coverage}
    \label{fig:image1}
  \end{subfigure}
  \begin{subfigure}{0.45\textwidth}
    \centering
    \includegraphics[width=\textwidth]{figures/prec_cov2.png}
    \caption{Precision Coverage}
    \label{fig:image2}
  \end{subfigure}
  
  \begin{subfigure}{0.45\textwidth}
    \centering
    \includegraphics[width=\textwidth]{figures/rec_ar2.png}
    \caption{Recall AMR}
    \label{fig:image3}
  \end{subfigure}
  \begin{subfigure}{0.45\textwidth}
    \centering
    \includegraphics[width=\textwidth]{figures/prec_ar2.png}
    \caption{Precision AMR}
    \label{fig:image4}
  \end{subfigure}
  
  \caption{GEOM-DRUGS evaluation as a function of number of generated conformers. GEOM-DRUGS has 104 conformers per molecule on average. Recall is heavily dependent on the sampling budget. Precision is mostly stable. Lower AMR and higher coverage is better, but coverage is set by an arbitrary threshold, which in this case is 0.75\AA.
  Results show \model (CC), Torsional Diffusion (TD), GeoMol (GM), and GeoDiff (GD).}
  \label{fig:2x2grid}
\end{figure}

\paragraph{GEOM-DRUGS Recall Results.} \fref{fig:2x2grid} demonstrates extensive Precision and Recall results for a wide range of tested sampling budgets for GEOM-DRUGS. We see that only Precision is stable across nearly all values. Due to the extreme sensitivity of the Recall metric and little difference in model performance for reasonable sampling budgets, we focus on Precision for QM9 and DRUGS. We also note that while \modelns-OT saw worse recall results for QM9, this was not the case for DRUGS. In the case of DRUGS, \modelns-OT achieves the learning objective of instilling force field optimizations as the lower error bound and does so with very little training and inference time.

% The RDKit+MMFF+Perfect Angles experiment was created to better understand where current ML and cheminformatic methods are making mistakes in conformer generation. This result demonstrates that the direct training objective of Torsional Diffusion and \modelns's torsion angle coarse-graining is on the right track, but more must be done to effectively learn the optimal torsion angles. We direct the reader to the stark difference between RDKit-MMFF and the Perfect Angle result. While \model is able to improve upon MMFF, it is a minor step compared to optimal angles. Surprisingly, although TD is trained to learn the optimal angle distribution, it is unable to do so. As a result, neither diffusion nor regression models have so far shown the ability to learn the optimal torsion angles. We leave future work to explore this further as we believe this is central to achieving the most accurate conformer generation.

\section{Oracle-based Protein Docking}
\label{section:protein}
We utilize the oracle-based protein docking task as molecules with higher affinity (more negative) have more potential for higher bioactivity, which is significant for real-world drug discovery. We use the CrossDocked2020 trainset consisting of 166000 protein-ligand interactions (2,358 unique proteins and 11,735 unique ligands) and its associated benchmarks, as it has been heavily used in Structure-based drug discovery as defined by~\cite{pocket2mol, targetdiff}.

The CrossDocked2020 dataset is derived from PDBBind but uses smina~\citep{smina}, a derivative of AutoDock Vina with more explicit scoring control, to generate the protein-conditioned ligand structures to yield ground truth data. We note that based on the raw data, 2.2 billion conformer-protein interactions are possible, but we filtered out any ground-truth example that AutoDock Vina failed to score. Furthermore, in the TDC oracle-based task, each ligand is known to fit well in the given protein. CrossDocked2020, on the other hand, consists of various ligand-protein interactions, not all of which are optimal, making the overall task more difficult.

We note that while it takes on the orders of hours to generate 1.2M conformers (100 conformers per molecule), it takes on the orders of \textasciitilde{weeks to months} to score each conformer for up to the 2,358 unique proteins for each evaluated method (evaluation time is 100x the time to score the ground truth data as we generate 100 conformers per molecule). As a result, we report the results for the first 100,000 conformer-protein interactions.

\section{Limitations}
\label{section:appendix-limits}
\vspace{-1ex}
As demonstrated in \sref{sec:rmsd}-\sref{sec:rigid}, \model significantly improves the accuracy and reduces
the overall data usage and runtime for conformer generation. However, \model also has some limitations that we will discuss in this section.
% \danny{skip for submission}
% \paragraph{MCG Benchmarks.} While MCG has been a prominent ML task for a couple of years, the primary RMSD benchmarks do not result in anything chemists or material scientists could use. They are mainly concerned with generating low-energy structures, and while the xTB property prediction task is a good start, the fact conformer energy is not used in the training process, nor is the model concerned with generating the lowest energy structure, is a major domain issue. We continue to work with domain scientists to bridge the gap between physical science and machine learning, as represented by our protein docking evaluations. We encourage future work to continue this focus of domain applications to build scientifically useful innovations.

\paragraph{Autoregressive generation.} While \model improves accuracy with reduced training time and overall data, autoregressive generation is the main bottleneck in inference time. We linearize the input molecule based on spatially significant subgraphs and then process each one autoregressively. For a model with k torsion angles, we need $k+1$ passes through our decoder. 
Coarse-graining is an effective strategy to reduce the number of decoder passes compared to traditional atom-wise autoregressive modeling. For example, for a given molecule, the number of torsion angles (which we use to coarse-grain) is significantly less than the number of atoms.
Our choice of coarse-graining strategy allows us to break the problem into more manageable subunits, making autoregressive modeling a useful strategy, as it provides greater flexibility and control by allowing conditional dependence.
\model is a good example of the trade-offs that exist between generative flexibility and speed. We target this limitation by introducing a non-autoregressive version with an optimal transport loss. We see this improves the overall GEOM results, at the slight cost of a higher right tail of the error distribution.

\paragraph{Optimal Transport.} While \modelns-OT trained with a non-autoregressive decoder with an optimal transport loss significantly outperforms prior methods and accomplishes the goal of effectively learning from traditional cheminformatics methods, the recall results still have room for improvement, especially for QM9.
While \model (no OT) achieves competitive results, we believe that continuing to focus on how to better integrate physics and cheminformatics into machine learning will be crucial for improving downstream performance. Due to the above concerns, we chose to evaluate our non-OT model on the property prediction and protein docking tasks, as wanted to use our best model denoted by the lowest overall RMSD error distribution.


\paragraph{Approximate structure error.} 
The success of learning the optimal distortion between low-energy and RDKit approximate structure depends on having reasonable approximations.
While \model relaxes the rigid local structure assumption of Torsional Diffusion in a way that leverages the torsional flexibility in molecular structures, it still depends on an approximate structure. This is a non-issue in some instances, as RDKit does well. In more experimental cases for larger systems, the RDKit errors may be too significant to overcome. 
We emphasize that the underlying framework of \model is adjustable and can learn from scratch, not only the distortion from approximate RDKit structures. In some cases, this may be more appropriate if the approximations have particularly high error. 
We leave to future work to explore the balance between the approximation error and the inductive bias of learning from approximate structures, as well as methods to maintain flexibility while avoiding the issues of conditioning on out-of-distribution poor approximations.

\paragraph{Equivariance.} As \model uses the EGNN~\citep{EGNN} framework as its equivariant backbone and thus only scalar and vector operations, there is no simple way to incorporate higher-order tensors. As the value of using higher-order tensors is still actively being explored, and in some cases, the costs outweigh the benefits, we used simple scalar and vector operations and avoided expensive tensor products.
We leave exploring the use of higher-order equivariant representations to future work, as it is still an ongoing research effort~\citep{equivariance-survey}.